\section*{Research Experience}
\textbf{SEAP (Solar, Earth and Planetary Physics) Research Group} -- University of Antioquia, Colombia
\begin{itemize}
    \item \textbf{Young Researcher Intern (Oct. 2024 -- Present)} — Lead developer of \href{https://github.com/seap-udea/PRisma}{\texttt{PRisma}}, a Bayesian pipeline that detects exoplanetary rings via \emph{asterodensity profiling}: comparing transit-inferred stellar densities against independent estimates to reveal the ``PhotoRing effect'' caused by unmodeled rings.
    \item Built a deterministic forward model of ringed-transit observables (\texttt{exorings}) and a KDE-based likelihood, running nested-sampling inference (\texttt{dynesty}) to recover posteriors and Bayesian evidence.
    \item Designed and ran parallelized nested-sampling campaigns on compute-cluster resources to process TTV photometry posteriors for the Kepler-51 system (planets b and d, the archetypal ``super-puff'' exoplanets). Lead author of the resulting manuscript (in preparation, 2026).
    \item Co-developer of \href{https://github.com/seap-udea/pryngles}{\texttt{Pryngles}} (PlanetaRY spaNGLES), an open-source Python package for visualizing ringed exoplanets and modeling their transit light curves and polarimetric signals; contributed the modern \texttt{System} interface and its \href{https://pryngles.readthedocs.io}{technical documentation}.
\end{itemize}
